\section{Spin}

A spinor is something that does not come back to itself after one full turn. Turn it once around and it is the negative
of what it was, and only a second full turn returns it. This sounds exotic, but on the mesh it is a plain fact of a
finite group.

The twenty-four directions of a dock form the binary tetrahedral group, the double cover of the rotation group of
three-dimensional space. In that group a full rotation composes to minus one and a double rotation to plus one, exactly,
by composition in the finite group. So the spinor sign is not imposed, it is arithmetic in the cell, and the Clifford
relations, the algebra that says how spinors transform, hold exactly on the substrate.

The consequence a physicist checks first comes out right with nothing put in. The magnetic strength of the resulting
spin-half particle, its g-factor, is two, the electron's value, a number measured to twelve digits, and here it falls
out of the structure rather than being set. Spin-statistics, that half-integer spins are fermions obeying exclusion,
holds on the same footing.

This is the geometric root of matter. The cell that tiles spacetime is the same object as the double cover that carries
spin, so the handedness of matter and the lattice of spacetime are one thing seen twice. Place tones on that cell and
read how they transform, and they transform as spinors, before any particle content is chosen. The next section turns
this into the actual quarks and leptons.
